\documentclass[11pt,reqno]{amsart}
\usepackage[T1]{fontenc}
\usepackage[utf8]{inputenc}
\usepackage{lmodern,amssymb,mathtools,booktabs}
\usepackage[margin=1in]{geometry}
\usepackage[colorlinks=true,linkcolor=blue,citecolor=blue,urlcolor=blue]{hyperref}
\newtheorem{theorem}{Theorem}[section]
\newtheorem{lemma}[theorem]{Lemma}
\newtheorem{corollary}[theorem]{Corollary}
\theoremstyle{remark}
\newtheorem{remark}[theorem]{Remark}
\newcommand{\Lag}{\mathcal{T}}
\newcommand{\cenLag}{\mathcal{H}}
\newcommand{\Dlt}{\Delta}
\newcommand{\ind}{\mathbf{1}}
\newcommand{\QQ}{\mathbb{Q}}
\newcommand{\NN}{\mathbb{N}}
\newcommand{\CC}{\mathbb{C}}
\title{Discrete Laguerre inequalities for a colored partition family}
\author{1200}
\date{30 October 2026; public-release preparation}
\subjclass[2020]{Primary 05A17, 05A20; Secondary 11P82}
\keywords{restricted partitions, discrete Laguerre inequality, quasi-polynomial, periodic coefficient}
\begin{document}
\begin{abstract}
For integers \(d\ge2\) and \(m\ge2\), let \(p_{d,m}(n)\) be the coefficient of
\((1-z)^{-2d}(1-z^m)^{-1}\). We determine exactly when all discrete Laguerre
inequalities of orders at most \(d\) for this sequence hold strictly for all
sufficiently large indices. This occurs for every \(m\) when \(d\) is even,
and precisely for \(m\le 1+\frac12\binom{2d}{d}\) when \(d\) is odd.
Above the latter threshold, the order-\(d\) transform is eventually negative
on a specified residue class. At equality, its leading term vanishes on that
class; we compute a positive term two degrees lower. The argument gives a
cyclic-difference formula for the boundary coefficient of the order-\(d\)
transform of exact quasi-polynomials whose coefficients above degree
\(\ell-2d\) are constant. Rational residue polynomials provide explicit,
parameter-dependent thresholds and independently checkable sign certificates.
\end{abstract}
\maketitle

\section{Introduction}

For a real sequence \(f=(f(n))_{n\ge0}\), define its discrete Laguerre transform
of order \(d\ge1\) by
\begin{equation}\label{eq:lag}
 \Lag_d(f)(n)=
 \sum_{j=0}^{2d}(-1)^{j+d}\binom{2d}{j}f(n+j)f(n+2d-j).
\end{equation}
The corresponding inequality is \(\Lag_d(f)(n)\ge0\). In order one this is
twice the usual adjacent log-concavity difference, with a shift of the index.
We distinguish non-strict inequalities from eventual strict positivity.

Gajdzica proved sufficient conditions under which quasi-polynomial-like
functions satisfy discrete Laguerre inequalities
\cite[Theorems 4.3--4.4]{Gajdzica_2024}.
His criterion requires the relevant upper coefficients to be constant.
Examples 4.5--4.6 of that article demonstrate that this condition is not
necessary; the latter example treats a multiset with four ones and one part
of size \(300\). Problem 5.6 asks for more effective criteria.
The present paper studies the first periodic coefficient and gives a complete
eventual classification for the family
\begin{equation}\label{eq:gf}
 \sum_{n\ge0}p_{d,m}(n)z^n=\frac{1}{(1-z)^{2d}(1-z^m)}.
\end{equation}
Here the \(2d\) parts of size one have distinct colors. The smooth asymptotic
coefficient from the earlier criterion and the isolated example just
mentioned are existing results. Our principal contribution is the
all-parameter classification, including the odd critical case.

\begin{theorem}\label{thm:main}
Let \(d\ge2\) and \(m\ge2\) be integers. There is an explicitly computable
integer \(N(d,m)\) such that
\[
 \Lag_j(p_{d,m})(n)>0
 \qquad(1\le j\le d,\ n\ge N(d,m))
\]
if and only if either \(d\) is even, or \(d\) is odd and
\begin{equation}\label{eq:cutoff}
 m\le 1+\frac12\binom{2d}{d}.
\end{equation}
If \(d\) is odd and \eqref{eq:cutoff} fails, then
\(\Lag_d(p_{d,m})(n)<0\) for every sufficiently large integer \(n\) satisfying
\(m\mid n+2d\). Explicit thresholds are supplied by the finite coefficient
formulas in Sections \ref{sec:cert} and \ref{sec:partition}.
\end{theorem}

The bound depends on the fixed parameters \(d,m\); no bound independent of
both parameters is asserted. A few consequences illustrate the parity
distinction:
\begin{center}
\begin{tabular}{ccc}
\toprule
Order \(d\) & Allowed \(m\) for eventual strict positivity & Critical value\\
\midrule
\(2\) & every \(m\ge2\) & none\\
\(3\) & \(2\le m\le11\) & \(11\)\\
\(4\) & every \(m\ge2\) & none\\
\(5\) & \(2\le m\le127\) & \(127\)\\
\bottomrule
\end{tabular}
\end{center}
These entries are consequences of the theorem, rather than a finite-search
claim.

We work with exact quasi-polynomials. This is a specified subclass of the
broader functions in the source problem: an arbitrary unspecified
little-\(o\) remainder does not provide an explicit error modulus.
We do not claim a classification for all restricted partition multisets
or for arbitrary periodic upper coefficients.

\section{The first periodic coefficient}\label{sec:boundary}

A quasi-polynomial of period \(M\ge1\) is specified by polynomials
\(f_0(X),\ldots,f_{M-1}(X)\), with \(f(n)=f_r(n)\) when \(n\equiv r\pmod M\).
For a periodic vector \(b\), all its arguments are understood modulo \(M\).
Put
\[
 \Dlt b(t)=2b(t)-b(t-1)-b(t+1).
\]
It is useful to center \eqref{eq:lag}:
\[
 \cenLag_d(f)(t)=\sum_{j=-d}^{d}w_jf(t+j)f(t-j),
 \qquad w_j=(-1)^j\binom{2d}{d+j}.
\]
Then \(\Lag_d(f)(n)=\cenLag_d(f)(n+d)\), and \(w_j=w_{-j}\).

\begin{lemma}\label{lem:moment}
For the preceding weights,
\[
 \sum_{j=-d}^{d}w_jj^k=0\quad(0\le k<2d),
 \qquad
 \sum_{j=-d}^{d}w_jj^{2d}=(-1)^d(2d)!.
\]
Moreover, \(\sum_jw_jb(t+j)=\Dlt^db(t)\).
\end{lemma}
\begin{proof}
The finite binomial identities
\[
 \sum_j w_jz^j=(2-z-z^{-1})^d,\qquad
 \sum_jw_je^{jv}=(2-2\cosh v)^d
\]
give the cyclic-shift formula and the required moments. Indeed, the second
expression has expansion \((-1)^dv^{2d}+O(v^{2d+2})\) as a formal power series.
\end{proof}

\begin{theorem}\label{thm:boundary}
Let \(d\ge1\), \(\ell\ge2d\), and let \(P\) be a real polynomial of degree
\(\ell\) and leading coefficient \(a>0\). Suppose
\[
 f_r(X)=P(X)+b_rX^{\ell-2d}+Q_r(X),\qquad
 \deg Q_r\le\ell-2d-1.
\]
When \(\ell=2d\), the last condition means \(Q_r=0\).
The residue polynomial of \(\Lag_d(f)\) has degree at most \(2\ell-2d\);
its coefficient at this degree on residue \(r\) is
\begin{equation}\label{eq:gamma}
 \gamma_r=(2d)!\binom{\ell}{d}a^2+2a\Dlt^db(r+d).
\end{equation}
For \(1\le j<d\), the residue polynomials of \(\Lag_j(f)\) have degree
\(2\ell-2j\) and positive, residue-independent leading coefficient
\((2j)!\binom{\ell}{j}a^2\).
\end{theorem}
\begin{proof}
The coefficient of \(t^k\) in \(P(t+j)P(t-j)\) has degree at most
\(2\ell-k\) as a polynomial in \(j\). Lemma \ref{lem:moment} therefore
annihilates the terms with \(k>2\ell-2d\). At equality, any term involving
a lower-degree monomial of \(P\) still has \(j\)-degree less than \(2d\).
The leading monomials give
\[
 a^2(t+j)^\ell(t-j)^\ell=a^2(t^2-j^2)^\ell.
\]
The coefficient of \(t^{2\ell-2d}\) before weighted summation is
\(a^2(-1)^d\binom{\ell}{d}j^{2d}\). This yields the first term of
\eqref{eq:gamma}.

Write \(R=f-P\). On a fixed residue of \(t\), the leading term of
\(R(t+j)\) is \(b(t+j)t^{\ell-2d}\). The cross terms contribute
\[
 a\sum_jw_j\bigl(b(t+j)+b(t-j)\bigr)=2a\Dlt^db(t)
\]
at degree \(2\ell-2d\). Products of two \(R\) terms have degree at most
\(2\ell-4d\), and terms involving \(Q_r\) contribute below \(2\ell-2d\).
The translation \(t=n+d\) preserves the highest possible coefficient and
gives the residue shift in \eqref{eq:gamma}.
For \(j<d\), the periodic term is below the boundary degree for order \(j\);
the same argument gives the last assertion.
\end{proof}

If all \(\gamma_r\) are positive, the order-\(d\) transform is eventually
strictly positive. A negative \(\gamma_r\) gives infinitely many failures.
A zero \(\gamma_r\) supplies no conclusion without lower coefficients.
The latter distinction is essential in the proof of Theorem \ref{thm:main}.

\section{Explicit bounds and residue certificates}\label{sec:cert}

For a nonzero polynomial \(h(X)=c_kX^k+\cdots+c_0\), define
\begin{equation}\label{eq:bound}
 B(h)=1+\left\lceil
       \frac{\sum_{i=0}^{k-1}|c_i|}{|c_k|}
       \right\rceil.
\end{equation}
The sum is empty for a constant polynomial.

\begin{lemma}\label{lem:bound}
For every real \(X\ge B(h)\), \(h(X)\) has the sign of \(c_k\).
\end{lemma}
\begin{proof}
The constant case is immediate. Otherwise put \(S=\sum_{i<k}|c_i|\).
Since \(X\ge1\), the absolute value of the lower-degree part is at most
\(SX^{k-1}\). But \(X>S/|c_k|\), so the leading monomial strictly dominates.
\end{proof}

For exact rational input, the transform on residue \(r\) is the explicitly
known rational polynomial
\begin{equation}\label{eq:residue}
 h_{d,r}(X)=\sum_{j=0}^{2d}(-1)^{j+d}\binom{2d}{j}
 f_{r+j}(X+j)f_{r+2d-j}(X+2d-j).
\end{equation}
The residue subscripts are reduced modulo \(M\). A zero polynomial has
sign zero throughout that residue; every nonzero polynomial has the eventual
sign of its leading nonzero coefficient. For one order a common bound is
\[
 N=\max\bigl(\{1\}\cup\{B(h_{d,r}):h_{d,r}\ne0\}\bigr).
\]
For orders \(1,\ldots,d\), take the same maximum over the order as well.
This convention gives \(N=1\) even if all output polynomials vanish.
Zero polynomials never count as strictly positive.

Thus an exact certificate consists of input coefficient vectors, the vectors
in \eqref{eq:residue}, their actual leading signs, and the integer bounds
\eqref{eq:bound}. Rational arithmetic checks all identities and inequalities.
This elementary polynomial-sign mechanism is used for effectivity; it is
not itself presented as a new general principle.

\section{Exact decomposition of the partition family}\label{sec:partition}

Fix \(q=2d\) and \(m\ge2\). From \eqref{eq:gf},
\begin{equation}\label{eq:direct}
 p_{d,m}(n)=\sum_{v=0}^{\lfloor n/m\rfloor}
               \binom{n-mv+q-1}{q-1}.
\end{equation}
Set \(u_0=1/m\), and recursively define
\begin{equation}\label{eq:u}
 u_h=-\frac1m\sum_{k=1}^{\min(h,m-1)}
             (-1)^k\binom{m}{k+1}u_{h-k}
       \qquad(1\le h\le q).
\end{equation}
Let
\begin{equation}\label{eq:P}
 P(X)=\sum_{h=0}^qu_h\binom{X+q-h}{q-h},
 \qquad
 b_r=\binom{r+q-1}{q-1}-P(r)\quad(0\le r<m).
\end{equation}
Binomial expressions in \(X\) denote their polynomial expansions.

\begin{lemma}\label{lem:partition}
For every \(n\ge0\), \(p_{d,m}(n)=P(n)+b(n)\), where \(b\) is the periodic
extension of the vector in \eqref{eq:P}. The leading coefficient of \(P\) is
\begin{equation}\label{eq:a}
 a=\frac1{mq!}.
\end{equation}
Furthermore,
\begin{equation}\label{eq:bdiff}
 \Dlt^db(t)=\frac{(-1)^d}{m}
       \bigl(m\ind_{\{m\mid t+d\}}-1\bigr).
\end{equation}
\end{lemma}
\begin{proof}
Write \(S_m(z)=1+z+\cdots+z^{m-1}\). Its expansion at \(z=1\) is
\[
 S_m(1-y)=\sum_{k=0}^{m-1}(-1)^k\binom{m}{k+1}y^k.
\]
Thus \eqref{eq:u} gives the first \(q+1\) coefficients of its reciprocal.
The principal part of
\((1-z)^{-q-1}/S_m(z)\) at \(z=1\) is
\(\sum_{h=0}^qu_h(1-z)^{-q-1+h}\). Its coefficients give \(P\) and
\eqref{eq:a}.

At every nontrivial \(m\)th root of unity \(\zeta\), the remaining pole is
simple. Multiplication by \(1-z/\zeta\) and passage to the limit gives the
partial-fraction term
\[
 \frac{1}{m(1-\zeta)^q}\frac{1}{1-z/\zeta}.
\]
These terms account for all poles. Both sides are proper rational functions,
so their difference has no polynomial remainder. Hence
\begin{equation}\label{eq:broot}
 b(t)=\frac1m\sum_{\substack{\zeta^m=1\\\zeta\ne1}}
                 \frac{\zeta^{-t}}{(1-\zeta)^q}.
\end{equation}
This is periodic. When \(0\le r<m\), only \(v=0\) occurs in \eqref{eq:direct},
which verifies the rational formula for \(b_r\) in \eqref{eq:P}.

The cyclic difference acts on \(\zeta^{-t}\) by the multiplier
\[
 2-\zeta-\zeta^{-1}=-\frac{(1-\zeta)^2}{\zeta}.
\]
Raising this to \(d\) cancels the denominator in \eqref{eq:broot}, since
\(q=2d\). The sum of all \(m\)th roots raised to an integer power is \(m\)
if the power is divisible by \(m\), and zero otherwise. Removing the root
one proves \eqref{eq:bdiff}.
\end{proof}

Combining Theorem \ref{thm:boundary} and Lemma \ref{lem:partition}, the
coefficient of \(n^{2d}\) in the forward transform is
\begin{equation}\label{eq:familygamma}
 \frac{\binom{2d}{d}
       +2(-1)^d\bigl(m\ind_{\{m\mid n+2d\}}-1\bigr)}
      {m^2(2d)!}.
\end{equation}
All lower-order transforms have positive leading coefficients.

\section{The critical odd case}\label{sec:critical}

For even \(d\ge2\), the two numerators in \eqref{eq:familygamma} are
\(\binom{2d}{d}-2\) and \(\binom{2d}{d}+2(m-1)\), both positive.
For odd \(d\), the only potentially nonpositive numerator is
\(\binom{2d}{d}-2(m-1)\), on \(m\mid n+2d\).
This settles every case except
\begin{equation}\label{eq:critical}
 d\ge3\text{ odd},\qquad C=\binom{2d}{d},\qquad m=1+C/2.
\end{equation}
The parameter is an integer because \(C=2\binom{2d-1}{d-1}\).

Put \(q=2d\), \(\sigma=(q+m)/2\), and \(x=t+\sigma\). The polynomial part
has the centered expansion
\begin{equation}\label{eq:centerP}
 P(t)=ax^q+cx^{q-2}+O(x^{q-4}),\qquad
 c=-\frac{a q(q-1)(q+m^2)}{24}.
\end{equation}
Here and below the \(O\) notation in polynomial identities means a polynomial
of at most the indicated degree; no effective estimate is hidden in it.
To verify \eqref{eq:centerP}, write \(P(t)\) as minus the residue at \(z=1\)
of the generating function multiplied by \(z^{-t-1}\), and substitute
\(z=e^{-v}\). Formal residue invariance gives
\[
 P(t)=\frac1m[v^q]e^{(t+\sigma)v}
       \left(\frac{v}{2\sinh(v/2)}\right)^q
       \frac{mv}{2\sinh(mv/2)}.
\]
The product after the exponential is even and equals
\(1-(q+m^2)v^2/24+O(v^4)\), proving \eqref{eq:centerP}.

On a residue with \(m\mid t+d\), put
\[
 B_k=\sum_{j=-d}^{d}w_jj^k b(t+j).
\]
We require its first three moments:
\begin{equation}\label{eq:B}
 B_0=-\frac{m-1}{m},\qquad B_1=0,\qquad
 B_2=\frac{m-1}{m}
       \left(-d^2+\frac{d(2d-1)(m+1)}6\right).
\end{equation}
The first formula follows from \eqref{eq:bdiff}.
For the others, set
\(S(z)=(2-z-z^{-1})^d=(-1)^d(1-z)^qz^{-d}\) and
\(D=z\,d/dz\). Direct differentiation gives
\begin{equation}\label{eq:derivatives}
 \frac{DS}{S}=-d\frac{1+z}{1-z},
 \qquad
 \frac{D^2S}{S}=d^2+(4d^2-2d)\frac{z}{(1-z)^2}.
\end{equation}
In \eqref{eq:broot}, the critical residue has \(\zeta^{-t}=\zeta^d\).
The first moment is therefore a constant multiple of
\(\sum_{\zeta\ne1}(1+\zeta)/(1-\zeta)=0\), by pairing inverse roots.
The self-inverse root \(-1\), if present, contributes zero.
For the second moment use
\begin{equation}\label{eq:rootsum}
 \sum_{\substack{\zeta^m=1\\\zeta\ne1}}\frac{\zeta}{(1-\zeta)^2}
       =-\frac{m^2-1}{12}.
\end{equation}
Indeed, the logarithmic derivative of \(S_m\) at \(1\) gives
\(\sum(1-\zeta)^{-1}=(m-1)/2\) and
\(\sum(1-\zeta)^{-2}=(m-1)(5-m)/12\); subtraction proves
\eqref{eq:rootsum}. Equations \eqref{eq:broot} and
\eqref{eq:derivatives} now give \eqref{eq:B}.
The operator \(D\) acting on \(z^{-j}\) introduces \(-j\), whereas \(D^2\)
introduces \(j^2\); this sign distinction leaves \(B_1=0\) unchanged.

\begin{lemma}\label{lem:beta}
Under \eqref{eq:critical}, on the residue class \(m\mid t+d\), the residue polynomial of
\(\cenLag_d(p_{d,m})(t)\) has degree exactly \(q-2\), with positive leading
coefficient
\begin{equation}\label{eq:beta}
 \beta=
 \frac{a^2q!C\,d(m^2-4d^2)}{12}
 +2a\binom q2 B_2+2cB_0.
\end{equation}
\end{lemma}
\begin{proof}
The degree-\(q\) coefficient is zero by \eqref{eq:critical}.
There is no degree-\(q-1\) coefficient: the smooth part is even in \(x\),
and the only possible periodic cross term at this degree is a multiple
of \(B_1=0\).

For the next coefficient, expand the moment identity one term further:
\[
 \sum_jw_je^{jv}=(-1)^dv^q
                   \left(1+\frac{d}{12}v^2+O(v^4)\right).
\]
The \(a^2\) contribution to \(x^{q-2}\) is
\(-a^2\binom q{d+1}(q+2)!d/12\). The mixed smooth contribution is
\[
 2acq!\left(\binom{q-2}{d}-\binom{q-2}{d-1}\right).
\]
It follows either by direct expansion or by the identity
\[
 (x+j)^q(x-j)^{q-2}+(x+j)^{q-2}(x-j)^q
       =2(x^2+j^2)(x^2-j^2)^{q-2}.
\]
Substitute \eqref{eq:centerP},
\(\binom q{d+1}=Cd/(d+1)\), and
\(\binom{q-2}{d-1}=Cd/[2(2d-1)]\).
The two contributions reduce to the first term of \eqref{eq:beta}.
All lower smooth terms are annihilated at this degree by
Lemma \ref{lem:moment}. The cross terms with \(b\) contribute the remaining
two terms in \eqref{eq:beta}; the purely periodic product has degree zero,
which is below \(q-2\ge4\).

Every displayed summand in \eqref{eq:beta} is positive.
For \(d\ge3\),
\[
 C\ge\binom{2d}{2}=d(2d-1),\qquad
 1+\frac{d(2d-1)}2-2d=\frac{(2d-1)(d-2)}2>0.
\]
Thus \(m>2d\). Also \(c<0\), \(B_0<0\), and
\[
 (2d-1)(m+1)>(2d-1)(2d+1)=4d^2-1>6d,
\]
so \(B_2>0\). This proves \(\beta>0\).
\end{proof}

\begin{proof}[Proof of Theorem \ref{thm:main}]
The noncritical cases follow from \eqref{eq:familygamma}, and
Lemma \ref{lem:beta} resolves equality in \eqref{eq:cutoff}.
Every lower order has a positive leading coefficient by
Theorem \ref{thm:boundary}. Translation from \(t\) to \(n\) preserves
each actual leading sign. Formulas \eqref{eq:u}--\eqref{eq:P} give the exact
input residue polynomials; \eqref{eq:residue} gives each transform.
Taking the maximum of \eqref{eq:bound} over the finitely many nonzero
polynomials and adjoining \(1\) gives the asserted \(N(d,m)\).
Above the cutoff, the negative leading coefficient on
\(n\equiv-2d\pmod m\) gives the stated repeated failure, with an explicit
bound from the same construction.
\end{proof}

\begin{remark}
The restriction \(d\ge2\) is essential for this strict classification.
For \(d=1,m=2\),
\(p_{1,2}(n)=\lfloor(n+2)^2/4\rfloor\) and
\(\Lag_1(p_{1,2})(2k)=0\),
\(\Lag_1(p_{1,2})(2k+1)=2(k+2)^2>0\).
For \(m\ge3\), \eqref{eq:familygamma} has a negative residue.
\end{remark}

\section{Computational and formal verification}

The mathematical proofs above are algebraic and do not depend on finite
experiments. The accompanying code generates rational residue polynomials,
their transformed coefficient vectors, and the explicit sign bounds.
The initial certificate set includes the critical pairs \((3,11)\) and
\((5,127)\), the pairs immediately above and below the latter cutoff,
and the previously known pair \((2,300)\).
Two separate model sessions reviewed the paper argument; their single
minor threshold-convention issue was repaired.

Independent Newton interpolation from direct partition counts, followed by
recurrence and initial-value checks, reproduces all coefficients, signs,
and bounds in 724 partition records and 14 boundary test polynomials.
A counterexample search covered 63 parameter pairs (all orders up to the
specified order) and 56 rational periodic examples with a fixed seed.
All 41,705 exact identity checks passed; point and bound checks are included
in that total. These finite checks do not prove the infinite-parameter theorem.

The Lean-verified component is the algebraic cancellation
\[
 \frac{(2-z-z^{-1})^d}{(1-z)^{2d}}
       =\frac{(-1)^d}{z^d}
 \qquad(z\in\CC\setminus\{0,1\},\ d\in\NN),
\]
used in \eqref{eq:bdiff}. Further checked statements establish the sign criterion
for the explicit numerators in \eqref{eq:familygamma}, positivity of the
expression \eqref{eq:beta}, the transformation formula \eqref{eq:residue}
for a supplied rational residue representation, and the rational version of
the coefficient bound \eqref{eq:bound}, including a common finite bound.
These components pass Lake builds with Lean 4.35.0 and Mathlib and independent
blind statement comparisons. The axiom audits report only \texttt{propext},
\texttt{Classical.choice}, and \texttt{Quot.sound}, without admitted proofs,
custom axioms, or \texttt{native\_decide}.
The original targets now have author-supplied unconditional Lean proofs:
\texttt{main\_classification\_proved} and
\texttt{above\_threshold\_failure\_proved}. The concrete Bernoulli residue
construction is for the same finite partition sum; not every displayed
partial-fraction identity has been formalized verbatim. Author checks report
a clean project rebuild, axiom audits and \texttt{leanchecker --fresh}.
Independent statement and formal audits under the repository workflow remain
pending. The general real-coefficient bound is not claimed formalized.

\section{Concluding remarks}

The centered polynomial and two Fourier moments resolve the vanishing
boundary coefficient for \eqref{eq:gf}. Finding similarly explicit critical
coefficients for other finite multisets is a natural further question.

The documented search found no equivalent full classification. Access to
some databases and a related final Acta Arithmetica article was incomplete;
global publication absence is not proved. The uninspected article is not
used in the argument, and the known smooth criterion and earlier example
are explicitly credited above.

\section*{Use of AI tools}
This research draft was developed with OpenAI Codex language models during
21--30 October 2026, including
formulation of candidate statements, proofs, computational programs, and
exposition. Separate model sessions performed paper review, a novelty audit,
and a blind comparison of the Lean statement. Computational and formal
verification statuses are stated above. These are AI-assisted checks and
do not replace human authorship or editorial review. No AI system is listed
as an author. Human author review and sign-off are required before public
release or submission.
Exact deployed model build identifiers were not recorded.

\section*{Data and code availability}
The local research package contains the original and independent-check
programs, rational coefficient certificates, SHA-256 manifests, review
records, and Lean sources. The planned public repository is \url{https://github.com/1200r3/math-research-results}; its release and independent-verification status are tracked separately. The package is intended to accompany any subsequently authorized
submission.

\bibliographystyle{plainurl}
\bibliography{refs}
\end{document}

